\documentclass{amsart}
\usepackage{graphicx} % Required for inserting images
\usepackage[dvipsnames]{xcolor}
\usepackage{amsmath}
\usepackage{amssymb}
\usepackage{stackrel}
\usepackage{amsthm}
\usepackage[style=alphabetic,maxnames=99,minalphanames=3,maxalphanames=4]{biblatex}
\usepackage{mathrsfs}
\usepackage{outlines}
\usepackage{tikz-cd}
\usepackage{todonotes}
\usepackage{lacromay}
\usepackage[hidelinks, colorlinks=true,linkcolor=MidnightBlue, citecolor=ForestGreen, urlcolor=Violet]{hyperref}
\AtBeginBibliography{\footnotesize}
\renewcommand*{\labelalphaothers}{$^{+}$}
\setcounter{tocdepth}{3}% to get subsubsections in toc

\let\oldtocsection=\tocsection

\let\oldtocsubsection=\tocsubsection

\let\oldtocsubsubsection=\tocsubsubsection

\renewcommand{\tocsection}[2]{\hspace{0em}\oldtocsection{#1}{#2}}
\renewcommand{\tocsubsection}[2]{\hspace{1em}\oldtocsubsection{#1}{#2}}
\renewcommand{\tocsubsubsection}[2]{\hspace{2em}\oldtocsubsubsection{#1}{#2}}

\title{Connectivity}
\author{Ajay Srinivasan}
\address{\scriptsize{Department of Mathematics, University of Southern California, Los Angeles, CA 90007}}
\email{avsriniv@usc.edu}
\date{\today}

\addbibresource{biblio.bib}

\let\fullref\autoref

\def\makeautorefname#1#2{\expandafter\def\csname#1autorefname\endcsname{#2}}
\def\makeautorefname#1#2{\expandafter\def\csname#1autorefname\endcsname{#2}}


\def\equationautorefname~#1\null{(#1)\null}
\makeautorefname{footnote}{footnote}%
\makeautorefname{item}{item}%
\makeautorefname{figure}{Figure}%
\makeautorefname{table}{Table}%
\makeautorefname{part}{Part}%
\makeautorefname{appendix}{Appendix}%
\makeautorefname{chapter}{Chapter}%
\makeautorefname{section}{Section}%
\makeautorefname{subsection}{Section}%
\makeautorefname{subsubsection}{Section}%
\makeautorefname{thm}{Theorem}%
\makeautorefname{sta}{Statement}%
\makeautorefname{cor}{Corollary}%
\makeautorefname{lem}{Lemma}%
\makeautorefname{prop}{Proposition}%
\makeautorefname{pro}{Property}
\makeautorefname{conj}{Conjecture}%
\makeautorefname{defn}{Definition}%
\makeautorefname{notn}{Notation}
\makeautorefname{notns}{Notations}
\makeautorefname{rem}{Remark}%
\makeautorefname{rems}{Remarks}%
\makeautorefname{quest}{Question}%
\makeautorefname{exmp}{Example}%
\makeautorefname{ax}{Axiom}%
\makeautorefname{claim}{Claim}%
\makeautorefname{assn}{Assumption}%
\makeautorefname{asses}{Assumptions}%
\makeautorefname{countcom}{Assumption}%
\makeautorefname{countcoms}{Assumptions}%
\makeautorefname{countmcom}{Assumption}%
\makeautorefname{con}{Construction}%
\makeautorefname{prob}{Problem}%
\makeautorefname{warn}{Warning}%
\makeautorefname{obs}{Observation}%
\makeautorefname{conv}{Convention}%

\newtheorem{thm}{Theorem}[section]
\newtheorem{sta}{Statement}[section]
\newtheorem{cor}{Corollary}[section]
\newtheorem{prop}{Proposition}[section]
\newtheorem{lem}{Lemma}[section]
\newtheorem{prob}{Problem}[section]
\newtheorem{conj}{Conjecture}[section]


\theoremstyle{definition}
\newtheorem{defn}{Definition}[section]
\newtheorem{ax}{Axiom}[section]
\newtheorem{con}{Construction}[section]
\newtheorem{exmp}{Example}[section]
\newtheorem{notn}{Notation}[section]
\newtheorem{notns}{Notations}[section]
\newtheorem{pro}{Property}[section]
\newtheorem{quest}[defn]{Question}
\newtheorem{rem}{Remark}[section]
\newtheorem{rems}{Remarks}[section]
\newtheorem{warn}{Warning}[section]
\newtheorem{sch}{Scholium}[section]
\newtheorem{obs}{Observation}[section]
\newtheorem{conv}{Convention}[section]
\newtheorem{assn}{Assumption}[section]
\newtheorem{asses}{Assumptions}[section]

\newtheorem*{asses:monad}{Assumptions \ref{asses:monad}}

\renewcommand{\qedsymbol}{$\blacksquare$}
\newcommand{\Cmon}{C}
\newcommand{\Cpre}{C^{\mathrm{pre}}}
\newcommand{\sLU}{\mathscr{L}_{U}}

\newcommand{\TG}{\sT_G}
\newcommand{\UG}{\ul{G\sU_{\ast}}}

\newcommand{\FdashG}{\sF[G\sT]}
\newcommand{\PIdashG}{\PI[G\sT]}

\newcommand{\FsubG}{\sF_G[\ul{G\sT}]}
\newcommand{\PIsubG}{\PI_G[\ul{G\sT}]}

\newcommand{\DdashG}{\sD[\ul{G\sT}]}
\newcommand{\DsubG}{\sD_G[\ul{G\sT}]}

\newcommand{\EdashG}{\sE[\ul{G\sT}]}
\newcommand{\EsubG}{\sE_G[\ul{G\sT}]}

\newcommand{\Dalg}{\bD\big[\PI[\ul{G\sT}]\big]}
\newcommand{\DGalg}{\bD_G\big[\PI_G[\ul{G\sT}]\big]}

\newcommand{\mb}[1]{\mathbf{#1}}

\newcommand{\mI}{\mathcal{I}}
\newcommand{\mJ}{\mathcal{J}}

\newcommand{\gen}{$\bF_\bullet$}
\newcommand{\cof}{$\bF_\bullet$}

\newcommand{\lbull}{^{\bullet}S^{\star}}
\newcommand{\lbullv}{^{\bullet}S^{V}}
\newcommand{\lbullov}{^{\bullet}S_0^{V}}
\newcommand{\lbullo}{^{\bullet}S_0^{\star}}
\newcommand{\lbullvw}{^{\bullet}S^{V\oplus W}}

\newcommand{\rbull}{S^{\star}{^{\bullet}}}
\newcommand{\rbullo}{S_0^{\star}{^{\bullet}}}
\newcommand{\rbulll}{S_1^{\star}{^{\bullet}}}
\newcommand{\rbullov}{{S_0^V}{^{\bullet}}}
\newcommand{\rbulllv}{{S_1^V}{^{\bullet}}}

\newcommand{\rbullv}{S^{V}{^{\bullet}}}
\newcommand{\rbullpv}{S^{V'}{^{\bullet}}}
\newcommand{\rbullw}{S^{W}{^{\bullet}}}
\newcommand{\rbullvw}{S^{V\oplus W}{^{\bullet}}}
%\newcommand{\nSv}{^n\!S^V}
%\newcommand{\nSvw}{^n\!S^{V\oplus W}}
%\newcommand{\mSv}{^m\!S^V}
%\newcommand{\sSv}{^s\!S^V}
%\newcommand{\phSv}{^{\ph_j}\!S^V}

\newcommand{\nSv}{^n\!(S^V)}
\newcommand{\nSvw}{^n\!(S^{V\oplus W})}
\newcommand{\mSv}{^m\!(S^V)}
\newcommand{\sSv}{^s\!(S^V)}
\newcommand{\phSv}{^{\ph_j}\!(S^V)}

%\newcommand{\Svn}{S^{V}\!{^{n}}}
%\newcommand{\Svm}{S^{V}\!{^{m}}}
%\newcommand{\Svs}{S^{V}\!{^{s}}}

\newcommand{\Svn}{(S^{V})^n}
\newcommand{\Svm}{(S^{V})^m}
\newcommand{\Svs}{(S^{V})^s}


%\newcommand{\inj}{\mathcal{I}}
\newcommand{\inj}{\vartheta}

\newcommand{\bi}{\mathbf{i}}
\newcommand{\bj}{\mathbf{j}}
%\newcommand{\bk}{\mathbf{k}}
\newcommand{\bm}{\mathbf{m}}
\newcommand{\bn}{\mathbf{n}}
\newcommand{\bp}{\mathbf{p}}
\newcommand{\bq}{\mathbf{q}}
\newcommand{\br}{\mathbf{r}}
\newcommand{\bs}{\mathbf{s}}

\newcommand{\OGop}{G\mathscr{O}^{op}[\sT]}
\newcommand{\Cp}{\mathbb{C}^{\mathrm{pre}}}


\begin{document}
\maketitle

%We say that $E \in \sS$ is connective if the counit $\epz_{\GA} \colon \SI^{\infty}_{\GA} \OM^{\infty}_{\GA} E \rtarr E$ is a weak homotopy equivalence. Equivalently if $B(\SI^{\infty},\GA,\OM^{\infty}E) \rtarr E$ is a weak homotopy equivalence. We note that $\SI^{\infty}$ takes values in connective spectra. By the two out of the three property and the triangle identity, this is the same thing as saying $\et_{\GA} \colon \OM^{\infty}_{\GA} \SI^{\infty}_{\GA} \overline{\OM^{\infty}_{\GA} E} \rtarr \overline{\OM^{\infty}_{\GA}E}$ is a weak homotopy equivalence. Since $B$ commutes with $\OM^{\infty}$ up to weak homotopy equivalence, this is equivalent to $B(\SI^{\infty},\GA,\OM^{\infty}E) \rtarr E$ being a weak homotopy equivalence. Now suppose $E = \SI^{\infty}X$, then $B(\SI^{\infty},\GA,\GA X)$ is naturally isomorphic to $\SI^{\infty}_{\GA} \overline{\GA X}$. We know that $\ze \colon \overline{\GA X} \rtarr \GA X$ is a homotopy equivalence, and therefore $\SI^{\infty}_{\GA} \overline{\GA X} \rtarr \SI^{\infty}_{\GA} \GA X$ is a weak homotopy equivalence. By the untwisting isomorphism $\SI^{\infty}_{\GA} \GA X \iso \SI^{\infty}X$. We have shown that $\SI^{\infty}$ takes values in $\sS_C$ defined in this manner.\\

%In the motivic case, this argument can be extended to colimits to show that every connective spectrum (in the very effective sense, generated by suspension spectra under colimits and extensions) is connective in the above sense.\\

%Actually, the story is much simpler. If $E$ is connective (in the sense that $\sS_C$ is the smallest subcategory of $\sS$ such that a map of connective spectra is a weak homotopy equivalence iff its image under $\OM^{\infty}$ is one), we know that $B(\SI^{\infty},\GA,\OM^{\infty}E)$ is also connective. Therefore $B(\SI^{\infty},\GA,\OM^{\infty}E) \rtarr E$ is a weak homotopy equivalence iff $\OM^{\infty} B(\SI^{\infty},\GA,\OM^{\infty} E) \rtarr \OM^{\infty}E$ is one. However, we know that $\OM^{\infty}$ and $B$ commute upto weak homotopy equivalence. Therefore this is the same thing as knowing that $\zeta \colon \overline{\OM^{\infty}E} \rtarr \OM^{\infty}E$ is a weak equivalence, which we certainly know is true. 

New definition for connectivity: we say that $E \in \sS$ is connective if the counit $\epz_{\GA} \colon \SI_{\GA} \OM_{\GA} E \rtarr E$ is a weak homotopy equivalence (WHE). By the two out of three property and the triangle identity as in Peter’s previous email, this definition equates to saying that $E$ is connective if $\et_{\GA} \colon \OM_{\GA}\SI_{\GA} \overline{\OM_{\GA} E} \rtarr \overline{\OM_{\GA} E}$ is a WHE. Since $B$ commutes with $\OM$ up to WHE, this is again equivalent to saying that $B(\SI, \GA, \OM E) \rtarr E$ must be a WHE. This is all taken out of Peter's email. Now to what we were talking about yesterday.\\

I think the proof that this notion of connectivity is equivalent to the older notion of connectivity must be formal (once we have the simplicial prelims) as follows. In one direction, we know a map of (old) connective objects is a WHE iff its image under $\OM$ is one. Suppose $E$ is (old) connective, then we know that $B(\SI, \GA, \OM E)$ is connective as well (since $\SI$ takes values in connective objects, and $B$ takes levelwise connective objects to connective objects). Therefore to show that $E$ is (new) connective, it suffices to show that the induced map $\OM B(\SI,\GA,\OM E) \rtarr \OM E$ is a WHE. Since $\OM$ and $B$ commute up to WHE, the fact that $\overline{\OM E} = B(\GA, \GA, \OM E) \rtarr \OM E$ is a homotopy equivalence completes this direction. In the other direction, we note that $\sS_C$ is in practice the largest full subcategory such that its morphisms are WHEs iff their images under $\OM$ are WHEs as well. Let $E$ and $E'$ be arbitrary objects in the (new) connective subcategory. We then know that $B(\SI,\GA,\OM E) \rtarr E$ is a WHE, and likewise that $B(\SI, \GA, \OM E') \rtarr E'$ is a WHE. Assume $\OM E \rtarr \OM E'$ is a WHE. Since $\SI$ and $\OM$ preserve weak equivalences, and $B$ preserves weak equivalences on restriction to good objects (at the least), (perhaps) we can say that the induced map $B(\SI,\GA,\OM E) \rtarr B(\SI,\GA,\OM E')$ is a WHE. An application of the two out of three property tells us that every morphism in the subcategory of (new) connective objects is a WHE iff its image under $\OM$ is one. \\

An aside: I doubt that it can be formally proven that $\SI$ takes values in (new) connective objects, without knowing that it already takes values in (old) connective objects. If it could be, then thanks to the above we would have a formal proof that $\SI$ takes values in (old) connective objects -- but that's an assumption we make, not a result we obtain. That said, here's maybe how the new definition could be presented:

\begin{defn}\label{defn:0.1}
We say that $E \in \sS$ is connective if the counit map $\epz_{\GA} \colon \SI_{\GA} \OM_{\GA} E \rtarr E$ is a weak homotopy equivalence. We denote the subcategory of connective objects $\sS_C$
\end{defn}

We require that $\SI$ takes values in $\sS_C$. In practice one might alternatively choose to define $\sS_C$ to be the largest full subcategory of $\sS$ such that its morphisms are WHEs iff their images under $\OM$ are WHEs as well.

\begin{prop}
The above two definitions are equivalent. Namely, $\sS_C$ as in Defn. \ref{defn:0.1} is equivalent to the largest subcategory of $\sS$ whose morphisms are WHEs iff their images under $\OM$ are WHEs as well.
\end{prop}
\end{document}